\chapter[序論（諸言：背景と目的）]{諸言}

\section{背景}

企業が投資家に提供する情報は、決算短信や適時開示、報道記事など多様であるが、その中でも有価証券報告書は、企業の経営実態や将来見通しに関する相対的に体系立った情報を網羅的に含む点で重要である。とりわけ「事業等のリスク」欄は、企業が認識する主要なリスク要因を言語で整理し、継続的に開示するセクションであり、投資家が企業の不確実性を評価する上で中心的な手がかりになり得る。一方で実務的には、当該セクションが前年と類似した表現の反復になりやすいこと（いわゆるコピペ）も指摘されており、開示の“量”だけではなく、“内容がどれだけ更新されているか”という質的側面を捉える必要性が高まっている。開示文が更新される局面、すなわち前年には見られなかったリスク要因の追加や、強調点の変化、具体化・詳細化が生じる局面では、投資家が受け取る情報の新規性が高い可能性がある。一般に投資家の注意には制約があるため、反復的な記述よりも、新規性を伴う記述の方が市場参加者の注目を集め、売買行動を誘発しやすいと考えられる。また、新たなリスク情報は将来のキャッシュフローや分散に関する不確実性の認識を変化させ得るため、短期的には株価の不確実性（予想外の価格変動）の増大として観察される可能性がある。他方で、新規な開示が不確実性を単に高めるとは限らず、リスクの所在を明確化することで投資家の理解が進み、反応が限定的となる可能性もある。したがって、リスク開示の新規性が市場の反応に与える影響は理論的に一意ではなく、実証的に検証すべき論点である。先行研究では、開示文の分量や語調、可読性、あるいは定型性（ボイラープレート性）に着目した分析が蓄積されてきた。しかし、日本企業の有価証券報告書を対象に、「事業等のリスク」欄の年度間の内容変化そのものを定量化し、その変化が開示日周辺の市場反応（例：投資家の注目度や不確実性）とどのように関係するかを検証した研究は依然として限定的である。本研究は、リスク記述をチャンク単位に分割し、前年記述との意味的類似度にもとづいて開示の新規性（変化度）を定量化した上で、提出日をイベント日とするイベントスタディ枠組みにより、短期の市場反応を検証する。これにより、リスク開示が「形式的に存在する情報」なのか、それとも「更新（新規性）を通じて市場参加者の行動や不確実性に影響する情報」なのかを明らかにすることを目的とする。

\section{関連研究}

本研究は、有価証券報告書におけるリスク開示研究と、開示テキストの更新（新規性）を捉えるテキスト分析、さらに提出日を基準としたイベントスタディによる市場反応分析を接続するものである。
以下では、（1）リスク開示の情報内容、（2）リスク開示“更新”の検出、（3）日本語有報研究、（4）LLM/埋め込みを用いたテキスト指標化、の観点から先行研究を整理し、本研究の位置づけを示す。なお、会計・ファイナンス領域におけるテキスト分析の体系的整理としては \cite{LoughranMcDonald2016} がある。

まず、リスク開示が企業の実態や市場のリスク認識と関連し得ることは、リスク語辞書や頻度分析を通じて示されてきた。初期研究として、訴訟リスクと任意開示の関係に着目し、開示文書における注意喚起的な表現（cautionary language）の役割を検討する枠組みが提示されている\cite{NelsonPritchard2007}。
その後、米国 10-K を対象に、リスク関連語の頻度やリスク開示の量的指標が、株式リスク（β やボラティリティ）や提出直後の市場反応（CAR）と結びつくことが報告され、リスク開示が市場参加者にとって情報価値を持つ可能性が示されている\cite{Campbell2014}。

一方で、リスク開示は「リスクを増やした」という悪材料の通知であるだけでなく、情報環境の改善として評価され得るため、開示内容の質（とりわけ具体性）に着目した研究も重要である。具体性の高いリスク開示ほど提出時 CAR が正となり、またアナリスト予想精度の改善と関連することが示されており、開示の“質”が市場反応を左右し得る点が示唆される\cite{HopeHuLu2016}。

次に、本研究が中心的に扱う「新規性（更新）」に関して、単なる分量や頻度ではなく、“前年から何が変わったか”を捉える必要性が指摘されてきた。四半期報告（10-Q）のリスク要因更新をテキストマッチングで検出し、更新企業で提出直後に株価下落が観察されること、ただし反応が一部遅れて生じ得ることを示す研究は、リスク開示の更新が市場に新情報を与えるという解釈と整合的である\cite{Filzen2015}。
また、リスク要因の追加・削除といった“変化”が不確実性指標（例：VRP 等）と結びつくことを示す研究は、リスク開示の変化そのものが投資家の将来リスク認識に影響し得ることを示唆する\cite{Lyle2023}。さらに、開示テキストの内容が投資家のリスク認識と結びつくことも報告されている\cite{KravetMuslu2013}。
加えて、開示テキストの構造的変化や論点の推移を、トピックモデル（例：Latent Dirichlet Allocation）により捉える研究も存在し、開示内容の「変化」を定量的に扱うという点で本研究の問題意識と接続する\cite{Dyer2017}。

日本の有価証券報告書「事業等のリスク」欄に関しても、テキスト分析に基づく蓄積が存在する。たとえば開示量（語数等）と情報の非対称性の低下（アナリスト予想精度の向上）との関連が示される一方、量の増大が内容の平板化を伴い得る点、前年比較や具体的記載の必要性が示唆されている\cite{Noda2016}。
さらに、日本語のリスク辞書（専門家辞書＋機械学習辞書）を用いた研究では、リスク関連語の出現確率が株価リスク指標（ボラティリティ）と正に相関することが示される一方、前年比の変化のみではリスク変動を十分説明できないことが報告されている\cite{SatoIkedaInoue2021}。この結果は、「頻度の前年差分」だけではリスク構造の変化を捉えきれない可能性、すなわち“意味的な更新”を捉える指標の必要性を示唆する。

以上を踏まえると、先行研究は（i）リスク開示が市場・リスク指標と関連し得ること、（ii）開示の質（具体性等）が市場反応を左右し得ること、（iii）開示の更新（追加・削除）が不確実性と結びつき得ることを示してきた。しかし、日本企業の有報「事業等のリスク」欄について、年度間の意味的更新を安定的に定量化し、提出日をイベント日として市場反応を多面的に検証する枠組みは依然として限定的である。
そこで本研究は、リスク記述をチャンク単位に分割し、前年チャンク群との最大類似度に基づく変化度（新規性）指標を構築したうえで、提出日前後の短期反応をイベントスタディで捉える。特に、価格反応（CAR）に加えて、投資家の注目度を表す出来高指標や、不確実性を表す異常収益の絶対値集計等を併用することで、リスク開示の新規性が「注目」「不確実性」「価格」のどこに強く現れるかを識別する点に本研究の特徴がある。

\section{研究目的・研究課題と仮説}

\subsection{研究目的}

本研究の目的は、有価証券報告書「事業等のリスク」欄における記述の新規性（前年からの内容更新）をテキスト類似度に基づいて定量化し、その新規性が提出日周辺の資本市場反応に与える影響を実証的に明らかにすることである。とくに、リスク開示の影響は価格反応（CAR）だけに表れるとは限らないため、本研究では、価格反応に加えて投資家の注目度（出来高）および不確実性の高まり（異常収益の絶対値集計等）を併用し、リスク開示の新規性が市場のどの側面に強く反映されるかを検証する。

\subsection{研究課題（Research Question）}

上記目的を踏まえ、本研究は以下の研究課題を設定する。

\begin{itemize}
  \item RQ1：有報「事業等のリスク」欄の新規性が高いほど、提出日周辺で投資家の注目度（出来高）は高まるか。
  \item RQ2：新規性が高いほど、提出日周辺で市場の不確実性（予想外の価格変動の大きさ）は高まるか。
  \item RQ3：新規性が高いほど、提出日周辺で株価はネガティブに反応するか（CAR は低下するか）。
\end{itemize}

\subsection{仮説}

リスク開示の更新は投資家にとって新情報であり、注意の配分や解釈のばらつきを通じて取引活発化や不確実性の上昇をもたらし得る。また、リスク要因の追加・強調は一般にネガティブ情報として受け止められやすい。したがって、本研究では以下の仮説を設定する。

\begin{itemize}
  \item H1（注目度）：リスク開示の新規性が高いほど、提出日周辺の異常出来高（注目度）は増加する。
  \item H2（不確実性）：リスク開示の新規性が高いほど、提出日周辺の不確実性指標（例：異常収益の絶対値集計）は増加する。
  \item H3（価格反応）：リスク開示の新規性が高いほど、提出日周辺のCAR は低下する。
\end{itemize}

\subsection{検証の概要}

本研究では、リスク記述をチャンク単位に分割し、前年のチャンク集合との最大類似度から変化度（新規性）を算出することで説明変数を構成する。アウトカムは、（i）市場調整した異常出来高（注目度）、（ii）市場モデルに基づく異常収益の絶対値集計（不確実性）、（iii）市場モデル CAR（価格反応）とし、提出日をイベント日とするイベントスタディおよび回帰分析により、上記仮説を検証する。

\section{本研究の新規性・貢献}

本研究の新規性・貢献は以下のとおりである。第一に、有価証券報告書「事業等のリスク」欄における開示内容について、単なる文量やリスク語頻度ではなく、前年からの内容更新（新規性）に焦点を当てる点に新規性がある。リスク開示は年次で反復されやすく、記述の定型性が高い場合には、開示“量”が増えても情報価値が増えない可能性がある。そのため、開示が実際にどの程度更新されているのかを定量化し、市場がその更新に反応するかを検証することは、リスク開示の有用性を評価する上で重要である。第二に、本研究は新規性の測定において、リスク記述をチャンク単位に分割し、前年記述との意味的類似度に基づく変化度指標を構築する。これにより、表現の揺れや言い換えが存在しても、内容としての近さ・遠さを一定程度考慮した形で「更新」を捉えることができ、従来の単純な頻度差分では捉えにくい意味的変化を扱える点に方法論的貢献がある。第三に、市場反応を価格反応（CAR）だけに限定せず、投資家の注目度（出来高）および不確実性（異常収益の絶対値集計等）を併用して多面的に検証する。リスク開示の更新は、直ちに価格に強く反映されるとは限らず、まず注意の喚起や解釈のばらつきとして現れる可能性がある。本研究の設計は、リスク開示新規性が市場のどの側面に表れやすいかを識別し、実証結果の解釈可能性を高める。以上により、本研究は、日本企業の有報「事業等のリスク」欄を対象に、開示の“更新（新規性）”という質的側面が市場参加者の反応とどのように結びつくかを明らかにし、開示実務および投資家の情報処理に関する理解に資する知見を提供する。

\section{論文の構成}

本論文の構成は次のとおりである。第 1 章では研究背景、先行研究、研究目的・仮説、ならびに本研究の貢献を示した。第 2 章ではデータ、変数定義、新規性指標の構築方法、およびイベントスタディと回帰分析の設計を説明する。第 3章では実証結果を提示し、第 4 章で結果の解釈、含意、限界点を考察する。最後に第 5 章で結論と今後の課題を述べる。
